\section{Results and Discussion}

\subsection{Experimental Scope}

The results in this section come from the synthetic SEBA-XAI benchmark and
should be read only within that scope. The workload contains generated
stations, officers, records, policy contexts, and access requests. The
benchmark is useful for testing whether the architecture is measurable and
whether different audit designs fail in different ways. It is not evidence of
performance on live CCTNS or ICJS data.

\subsection{Baseline Access-Control and Audit Comparison}

Table~\ref{tab:method-comparison} shows the main prototype comparison on the
1,000-request paper evidence workload. A role/action-only RBAC baseline reaches
0.2900 accuracy against the declared policy oracle and creates 656 false
allows. This supports the design choice that role-only access is too weak for
the synthetic sensitive-record workload. The configured ABAC/PBAC policy
matches the declared oracle in this benchmark. However, the mutable-log version
does not provide internal tamper evidence. Adding a signed hash chain or
blockchain-style audit gives tamper evidence for controlled ordinary edits,
while the full SEBA-XAI mode additionally logs the explanation hash.

\begin{table}[t]
\centering
\caption{Selected method comparison on the 1,000-request evidence workload}
\label{tab:method-comparison}
\begin{tabular}{lcccc}
\toprule
Method & Acc. & Tamper & XAI hash & p50 ms \\
\midrule
RBAC + mutable log & 0.2900 & 0.0000 & No & 7.54 \\
ABAC/PBAC + mutable log & 1.0000 & 0.0000 & No & 20.99 \\
ABAC/PBAC + signed chain & 1.0000 & 1.0000 & No & 30.99 \\
ABAC/PBAC + blockchain audit & 1.0000 & 1.0000 & No & 42.09 \\
SEBA-XAI full & 1.0000 & 1.0000 & Yes & 42.09 \\
\bottomrule
\end{tabular}
\end{table}

The interpretation is not that blockchain is always better than a signed hash
chain. The local result only shows that both integrity-preserving append-only
designs detect the controlled ordinary tamper cases. The reason to keep the
blockchain-style layer in the architecture is the inter-agency audit setting:
known agencies can maintain a shared audit commitment without placing raw
records on-chain.

\subsection{Compromised-Signer Result}

The strongest technical result is the compromised-signer asymmetry. In the
five-seed adversarial grid, mutable logs, signed chains, blockchain-style
audit, CT-style logs, ABAC re-execution, and Fabric-style ABAC all record a
0.0000 detection rate for the validly re-signed compromised-signer attack.
NS-PI drift and the trusted raw-attribute policy oracle record a 1.0000
detection rate for that attack. This result does not mean NS-PI is the best
overall defense. Its mean adversarial audit score is 0.2500, while the signed
chain, blockchain-style audit, CT-style log, ABAC re-execution, and
Fabric-style ABAC each score 0.7917 across the broader attack set. The
trusted policy oracle scores 1.0000 because it has the strongest visibility
assumption: access to an uncompromised raw request view.

\begin{table}[t]
\centering
\caption{Five-seed audit results for selected defenses}
\label{tab:compromised-signer}
\begin{tabular}{lcc}
\toprule
Defense & AAS mean & Compromised signer \\
\midrule
Mutable log & 0.5000 & 0.0000 \\
Signed hash chain & 0.7917 & 0.0000 \\
Blockchain-style audit & 0.7917 & 0.0000 \\
CT-style log & 0.7917 & 0.0000 \\
ABAC/Fabric-style re-execution & 0.7917 & 0.0000 \\
NS-PI drift & 0.2500 & 1.0000 \\
Trusted policy oracle & 1.0000 & 1.0000 \\
\bottomrule
\end{tabular}
\end{table}

This is the main reason the paper treats the system as layered. Integrity
mechanisms answer whether the recorded log was changed. Policy re-evaluation
answers whether the recorded decision agrees with policy under the available
attributes. NS-PI answers a narrower log-only question: whether the signed
decision trace has shifted in a suspicious way. These are different questions,
so the defenses should not be treated as interchangeable.

\subsection{NS-PI Sensitivity and Workload Stress}

The NS-PI results are useful but limited. In global compromised-signer
sensitivity tests, NS-PI misses 2\% and 5\% corruption but detects the attack
from 10\% corruption onward in the current workload. For targeted attacks,
global drift can miss localized corruption because the total changed fraction
of the whole workload is small. Grouped station/district drift is more useful:
station-level grouped NS-PI reaches full detection from 50\% of targeted
eligible station rows, and district-level grouped NS-PI reaches full detection
from 25\% of targeted eligible district rows. Very small 10\% targeted
corruption is still missed.

Workload stress tests over 500, 1,000, 2,500, and 5,000 generated requests
show the same basic compromised-signer pattern at 25\% corruption: signed
chains remain at 0.0000 detection, while NS-PI global detection and the trusted
oracle remain at 1.0000. At 10\% corruption, the signal is weaker for smaller
workloads; for example, the 500-request baseline arm has 0.2000 NS-PI global
detection across five seeds, while the 1,000-request and larger baseline arms
reach 1.0000. This is a useful boundary condition for the paper: NS-PI needs
enough traffic and enough corruption signal to be reliable.

\subsection{XAI and Audit Reconstruction}

The XAI component is strongest as a structured trace. Across five seeds, trace
completeness is 1.0000, counterfactual coverage is 1.0000, counterfactual
validity is 0.9964, stable decision/reason rate is 1.0000, and audit
reconstruction rate is 1.0000. These results support the claim that the
prototype can produce reviewable decision artifacts and link them back to audit
events.

The natural-language explanation is weaker than the structured trace. The mean
decisive-attribute text coverage is 0.9310, but the full text coverage rate is
0.7810. Therefore the paper should not claim that the explanation text is
complete. A fair claim is that the structured XAI artifact is complete in this
benchmark, while the text renderer still needs improvement before stronger
human-facing explanation claims are made.

\subsection{Metadata, Latency, and Storage}

The metadata experiment compares a full metadata ledger with a minimized
commitment ledger. The full metadata ledger has a schema-level exposure score
of 1.0000 because fields such as decision, purpose, requester station, and
target station remain visible. The minimized commitment ledger has a score of
0.0000 in the same proxy because sensitive context is represented by hashes or
commitments. This is not a formal privacy proof; it only shows that schema
minimization reduces direct metadata exposure in the prototype.

The local latency/storage measurements are also modest-scope prototype
measurements. Policy-oracle and XAI generation takes 14.46 ms p50 for 1,000
requests in the benchmark artifact. Mutable-log build time is 6.84 ms p50,
signed-chain build and verification are 16.84 ms and 9.16 ms p50, and
permissioned blockchain-style build and verification are 11.10 ms and 2.50 ms
p50 in the local simulation. Storage per event is 459.82 bytes for the mutable
log, 756.86 bytes for the signed hash chain, and 353.50 bytes for the
blockchain-style audit artifact used in this measurement. These numbers should
not be interpreted as production throughput, but they show that the prototype
records latency and storage overhead in a reproducible way.
